% ============================================================================
% LANGENTWURF LE-06b: ir + estar (Bewegung + Lage)
% Spanisch Klasse 7 - Sequenz 6: Mi ciudad
% ============================================================================
% Tier: 1 (vollständig)
% Doppelstunde: 3-4 von 8
% Fachfarbe: #c41e3a (Karminrot)
% ============================================================================

\documentclass[11pt,a4paper]{article}

\usepackage[utf8]{inputenc}
\usepackage[T1]{fontenc}
\usepackage[ngerman]{babel}
\usepackage{geometry}
\usepackage{fancyhdr}
\usepackage{lastpage}
\usepackage{xcolor}
\usepackage{tcolorbox}
\usepackage{enumitem}
\usepackage{tabularx}
\usepackage{longtable}
\usepackage{array}
\usepackage{booktabs}
\usepackage{amssymb}

% ============================================================================
% KONFIGURATION
% ============================================================================

\newcommand{\fach}{Spanisch}
\newcommand{\fachfarbe}{c41e3a}
\newcommand{\jahrgangsstufe}{7}
\newcommand{\schulart}{Gesamtschule}
\newcommand{\stundenthema}{ir + estar (Bewegung und Lage)}
\newcommand{\reihenthema}{Mi ciudad}
\newcommand{\stundennummer}{06b}
\newcommand{\reihenstunden}{4}
\newcommand{\gerniveau}{A1}
\newcommand{\lehrwerk}{Qué pasa 1}
\newcommand{\lehrwerkseite}{Unidad 6}
\newcommand{\lerngruppengroesse}{24}

% ============================================================================
% LAYOUT
% ============================================================================
\geometry{left=2.5cm, right=2.5cm, top=2.5cm, bottom=2.5cm}
\definecolor{fach}{HTML}{\fachfarbe}
\definecolor{gray}{HTML}{666666}
\definecolor{lightgray}{HTML}{f5f5f5}
\definecolor{successgreen}{HTML}{2a7a4b}
\definecolor{warningorange}{HTML}{b35c00}
\definecolor{basis}{HTML}{2a7a4b}
\definecolor{standard}{HTML}{2c5aa0}
\definecolor{erweiterung}{HTML}{7b1fa2}

\tcbuselibrary{skins,breakable}

\newtcolorbox{infobox}[1][]{
    colback=lightgray,
    colframe=fach,
    fonttitle=\bfseries,
    title=#1,
    breakable
}

\newtcolorbox{scorebox}{
    colback=white,
    colframe=fach,
    fonttitle=\bfseries,
    title=Didaktik-Score,
    breakable
}

\pagestyle{fancy}
\fancyhf{}
\fancyhead[L]{\small\textcolor{gray}{\fach{} -- Jg. \jahrgangsstufe{} -- \stundenthema}}
\fancyhead[R]{\small\textcolor{gray}{LE-\stundennummer{} | Tier 1}}
\fancyfoot[C]{\small\thepage{} / \pageref{LastPage}}
\renewcommand{\headrulewidth}{0.4pt}

% Differenzierungsmarker
\newcommand{\B}{\textcolor{basis}{\textbf{[B]}}}
\newcommand{\Std}{\textcolor{standard}{\textbf{[S]}}}
\newcommand{\E}{\textcolor{erweiterung}{\textbf{[E]}}}

% ============================================================================
% DOKUMENT
% ============================================================================
\begin{document}

% ============================================================================
% DECKBLATT
% ============================================================================
\begin{titlepage}
    \centering
    \vspace*{2cm}
    {\large\textcolor{gray}{\schulart}}\\[2cm]
    {\Huge\textcolor{fach}{\textbf{Langentwurf}}}\\[0.5cm]
    {\large LE-\stundennummer{} | Tier 1 (vollständig)}\\[2cm]
    {\LARGE\textbf{\stundenthema}}\\[0.5cm]
    {\large Sequenz: \reihenthema}\\[2cm]
    \begin{tabular}{rl}
        \textbf{Fach:} & \fach{} (\gerniveau) \\[0.3cm]
        \textbf{Klasse:} & \jahrgangsstufe \\[0.3cm]
        \textbf{Lehrwerk:} & \lehrwerk, \lehrwerkseite \\[0.3cm]
        \textbf{Doppelstunde:} & \stundennummer{} (Std. 3-4 von 8) \\[0.3cm]
    \end{tabular}
    \vfill
\end{titlepage}

\tableofcontents
\newpage

% ============================================================================
% 1. BEDINGUNGSANALYSE
% ============================================================================
\section{Bedingungsanalyse}

\subsection{Lerngruppe}
\begin{tabular}{ll}
    Kursgröße: & \lerngruppengroesse{} SuS \\
    Jahrgangsstufe: & \jahrgangsstufe \\
    Sprachniveau: & \gerniveau{} (GER) -- Anfänger \\
    Fremdsprache: & 2. Fremdsprache (nach Englisch) \\
\end{tabular}

\subsection{Vorwissen aus 06a}
\begin{itemize}
    \item Orte in der Stadt: el cine, la plaza, el parque, la biblioteca, el supermercado, etc.
    \item Präpositionen: en, cerca de, al lado de
    \item Bestimmte Artikel: el, la, los, las
    \item Grundlegende Satzstrukturen mit ser und estar (Einführung)
\end{itemize}

\subsection{Differenzierung (intern)}
\begin{tabular}{|l|p{3.5cm}|p{3.5cm}|p{3.5cm}|}
\hline
& \B{} \textbf{Basis} & \Std{} \textbf{Standard} & \E{} \textbf{Erweiterung} \\
\hline
\textbf{Ziel} & BR: Rezeption & MR: Produktion mit Hilfe & GY: Freie Produktion \\
\hline
\textbf{ir} & Lückentext mit Vorgabe & Konjugation selbst & + Sätze mit Wegbeschreibung \\
\hline
\textbf{estar} & Ortsangaben zuordnen & Sätze bilden & Stadtplan beschreiben \\
\hline
\end{tabular}

\vspace{0.5em}
\textbf{WICHTIG:} Diese Marker sind NUR für die Lehrkraft!

% ============================================================================
% 2. SACHANALYSE
% ============================================================================
\section{Sachanalyse}

\subsection{Sprachliche Strukturen}

\textbf{Verb \textit{ir} (gehen/fahren):}
\begin{center}
\begin{tabular}{|l|l|l|}
\hline
\textbf{Person} & \textbf{Pronomen} & \textbf{ir} \\
\hline
1. Sg. & yo & voy \\
2. Sg. & tú & vas \\
3. Sg. & él/ella/usted & va \\
1. Pl. & nosotros/-as & vamos \\
2. Pl. & vosotros/-as & vais \\
3. Pl. & ellos/ellas/ustedes & van \\
\hline
\end{tabular}
\end{center}

\textbf{Struktur:} \textit{ir + a + Ort}
\begin{itemize}
    \item Voy \textbf{al} cine. (a + el = al)
    \item Vamos \textbf{a la} plaza.
    \item Van \textbf{al} parque.
\end{itemize}

\vspace{1em}

\textbf{Verb \textit{estar} (sein/sich befinden):}
\begin{center}
\begin{tabular}{|l|l|l|}
\hline
\textbf{Person} & \textbf{Pronomen} & \textbf{estar} \\
\hline
1. Sg. & yo & estoy \\
2. Sg. & tú & estás \\
3. Sg. & él/ella/usted & está \\
1. Pl. & nosotros/-as & estamos \\
2. Pl. & vosotros/-as & estáis \\
3. Pl. & ellos/ellas/ustedes & están \\
\hline
\end{tabular}
\end{center}

\textbf{Verwendung:} Ortsangaben und Lage
\begin{itemize}
    \item El cine \textbf{está} en la plaza.
    \item ¿Dónde \textbf{estás}?
    \item Estamos en el parque.
\end{itemize}

\subsection{Kontrastive Betrachtung}

\begin{tabular}{|l|p{5cm}|p{5cm}|}
\hline
& \textbf{Deutsch} & \textbf{Spanisch} \\
\hline
Bewegung & gehen zu, fahren zu & ir + a + Ort \\
\hline
Kontraktion & --- & a + el = al \\
\hline
Lage & sein (+ Präposition) & estar + en/cerca de/... \\
\hline
ser vs. estar & nur \textit{sein} & estar für Ort, ser für Identität \\
\hline
\end{tabular}

\textbf{Typische Interferenz:}
\begin{itemize}
    \item *``Yo voy a el cine'' statt ``Voy al cine'' (Kontraktion vergessen)
    \item *``El cine es en la plaza'' statt ``está'' (ser/estar-Verwechslung)
\end{itemize}

\subsection{Kulturelle Aspekte}
\begin{itemize}
    \item Spanische Städte: Plaza als zentraler Treffpunkt
    \item Paseo: Abendlicher Spaziergang als soziales Ereignis
    \item Unterschied Großstadt/Dorf bei Freizeitaktivitäten
\end{itemize}

% ============================================================================
% 3. DIDAKTISCHE ANALYSE
% ============================================================================
\section{Didaktische Analyse}

\subsection{Stellung in der Sequenz}
\begin{tabular}{|c|l|l|}
\hline
\textbf{Block} & \textbf{Thema} & \textbf{Status} \\
\hline
6a & Orte in der Stadt (Wortschatz) & abgeschlossen \\
\textbf{6b} & \textbf{ir + estar (Bewegung/Lage)} & \textbf{diese Stunde} \\
6c & Wegbeschreibung & folgt \\
6d & Mi ciudad + Sequenztest & folgt \\
\hline
\end{tabular}

\subsection{Kompetenzziele (Rahmenplan MV)}
\begin{itemize}
    \item \textbf{Kommunikativ:} Sprechen (an Gesprächen teilnehmen), Schreiben
    \item \textbf{Sprachlich:} Grammatik (ir, estar), Wortschatz (Orte)
    \item \textbf{Interkulturell:} Stadtleben in Spanien
    \item \textbf{Methodisch:} Entdeckendes Lernen, Partnerarbeit
\end{itemize}

\subsection{Legitimation}
\textbf{Rahmenplan Spanisch MV (Kl. 7-10):}
\begin{itemize}
    \item Kompetenzbereich: Kommunikative Kompetenz -- Sprechen/Schreiben
    \item Standard: ``Die SuS können sagen, wohin sie gehen und wo sich Orte befinden.''
\end{itemize}

% ============================================================================
% 4. STUNDENZIELE
% ============================================================================
\section{Stundenziele}

\subsection{Grobziel}
Die SuS erweitern ihre \textbf{kommunikative Kompetenz im Bereich Sprechen}, indem sie mit \textit{ir + a} ausdrücken, wohin sie gehen, und mit \textit{estar} beschreiben, wo sich Orte befinden.

\subsection{Feinziele (operationalisiert)}
\begin{tabular}{|c|p{7.5cm}|c|c|}
\hline
\textbf{Nr.} & \textbf{Die SuS können...} & \textbf{AFB} & \textbf{Diff.} \\
\hline
FZ1 & die Formen von \textit{ir} und \textit{estar} \textbf{nennen} und korrekt aussprechen & I & alle \\
\hline
FZ2 & mit \textit{ir + a + Ort} sagen, wohin sie gehen (inkl. Kontraktion \textit{al}) & II & alle \\
\hline
FZ3 & mit \textit{estar + Präposition} beschreiben, wo sich ein Ort befindet & II & \Std{}/\E{} \\
\hline
FZ4 & einen kurzen Text über ihre Stadt/ihren Stadtteil \textbf{verfassen} & III & \E{} \\
\hline
\end{tabular}

\vspace{1em}
\textbf{AFB-Verteilung:} AFB I: 25\% | AFB II: 50\% | AFB III: 25\%

% ============================================================================
% 5. METHODISCHE ANALYSE
% ============================================================================
\section{Methodische Analyse}

\subsection{Einstieg}
\begin{itemize}
    \item \textbf{Methode:} Stadtplan-Impuls: ``¿Adónde vas?'' -- ``¿Dónde está...?''
    \item \textbf{Begründung:} Aktivierung Vorwissen (Orte), Problemstellung
    \item \textbf{Alternative:} Kurzvideo: Jugendlicher beschreibt seinen Weg
\end{itemize}

\subsection{Erarbeitung}
\begin{itemize}
    \item \textbf{Methode:} Entdeckendes Lernen -- ir/estar-Formen aus Beispielen ableiten
    \item \textbf{Begründung:} Aktive Konstruktion statt passiver Rezeption
    \item \textbf{Scaffolding:} Farbcodierung (ir = blau, estar = grün), Tabelle schrittweise füllen
\end{itemize}

\subsection{Übungsphasen}
\begin{itemize}
    \item \textbf{Methode:} Chorisches Sprechen $\rightarrow$ Partnerarbeit $\rightarrow$ Stadtplan-Dialog
    \item \textbf{Progression:} Gelenkt $\rightarrow$ Halbfrei $\rightarrow$ Frei
    \item \textbf{Differenzierung:} Stadtplankarten mit/ohne Hilfen
\end{itemize}

\subsection{Medien}
\begin{itemize}
    \item Tafel/Whiteboard: Konjugationstabellen ir/estar
    \item Lehrwerk: \lehrwerk, \lehrwerkseite
    \item Arbeitsblatt: AB-06b.pdf
    \item Lernpfad: LP-06b.html (Tablets, optional)
    \item Stadtplan (Kopiervorlage)
\end{itemize}

% ============================================================================
% 6. VERLAUFSPLANUNG
% ============================================================================
\section{Verlaufsplanung}

\textbf{1. Stunde (45 Min.)}

\begin{longtable}{|p{1cm}|p{1.5cm}|p{4cm}|p{3cm}|p{1.5cm}|p{2cm}|}
\hline
\textbf{Zeit} & \textbf{Phase} & \textbf{Lehrerhandeln} & \textbf{Schülerhandeln} & \textbf{SF} & \textbf{Material} \\
\hline
\endhead

0--5 & Einstieg & Stadtplan zeigen: ``¿Adónde vas los sábados?'' & Vermutungen, Orte nennen & Plenum & Stadtplan \\
\hline

5--10 & Aktivier. & Orte wiederholen, Beispielsätze: ``Voy al cine'' & Nachsprechen, Muster erkennen & Plenum & Tafel \\
\hline

10--20 & Input ir & Verb \textit{ir} einführen: Formen an Tafel, Kontraktion \textit{al} & Tabelle ergänzen, Beispiele bilden & Plenum & Tafel, AB \\
\hline

20--30 & Übung 1 & Chorisches Sprechen: ``Voy, vas, va...'' + Sätze & Nachsprechen, in PA üben & Plenum/PA & -- \\
\hline

30--40 & Übung 2 & PA: Minidialoge ``¿Adónde vas?'' -- ``Voy a...'' & Dialoge mit Stadtplan & PA & Stadtplan \\
\hline

40--45 & Sicherung & 2-3 Paare präsentieren, Merkkasten ir ausfüllen & Präsentieren, AB ergänzen & Plenum & AB-06b \\
\hline
\end{longtable}

\vspace{1em}

\textbf{2. Stunde (45 Min.)}

\begin{longtable}{|p{1cm}|p{1.5cm}|p{4cm}|p{3cm}|p{1.5cm}|p{2cm}|}
\hline
\textbf{Zeit} & \textbf{Phase} & \textbf{Lehrerhandeln} & \textbf{Schülerhandeln} & \textbf{SF} & \textbf{Material} \\
\hline
\endhead

0--5 & Aktivier. & Schnelles Quiz: ``Wie heißt \textit{ich gehe}?'' & Antworten: voy & Plenum & -- \\
\hline

5--15 & Input estar & Verb \textit{estar} einführen: ``El cine está en la plaza'' & Formen ableiten, Tabelle & Plenum & Tafel, AB \\
\hline

15--25 & Übung 3 & Stadtplan: ``¿Dónde está el cine?'' -- ``Está en...'' & Orte beschreiben & PA & Stadtplan \\
\hline

25--35 & Anwendung & Kombinationsübung: ir + estar im Dialog & Komplexere Dialoge & PA/GA & AB-06b \\
\hline

35--40 & Sicherung & Merkkasten estar ausfüllen, Unterschied ir/estar & AB ergänzen & Plenum & AB-06b \\
\hline

40--45 & Abschluss & Zusammenfassung, HA: Lehrwerk + LP-06b & Notieren & Plenum & -- \\
\hline
\end{longtable}

\textbf{Legende:} EA = Einzelarbeit, PA = Partnerarbeit, GA = Gruppenarbeit

% ============================================================================
% 7. ERWARTETE SCHWIERIGKEITEN
% ============================================================================
\section{Erwartete Schwierigkeiten}

\begin{tabular}{|p{4.5cm}|p{8cm}|}
\hline
\textbf{Schwierigkeit} & \textbf{Reaktion} \\
\hline
Kontraktion \textit{al} vergessen & Regelvisualisierung: a + el = al (Tafel), Kontrastübung \\
\hline
Verwechslung ir/estar & Farbcodierung: ir (Bewegung) = blau, estar (Lage) = grün \\
\hline
ser/estar-Verwechslung bei Ort & Merksatz: ``estar = wo, ser = wer/was'' \\
\hline
Akzente bei estás, está & Betonung üben, Bedeutungsunterschied esta/está \\
\hline
Zeitproblem & Puffer: Aufgabe 4 als HA, Kulturvergleich kürzen \\
\hline
\end{tabular}

% ============================================================================
% 8. HAUSAUFGABE
% ============================================================================
\section{Hausaufgabe}

\begin{itemize}
    \item Lehrwerk Unidad 6, Übung zu ir/estar
    \item Vokabeln lernen: ir-Formen, estar-Formen, Orte
    \item Optional: Lernpfad LP-06b.html bearbeiten
    \item \E{}: Kurzer Text ``Mi ciudad'' (5 Sätze mit ir/estar)
\end{itemize}

% ============================================================================
% 9. LÖSUNGEN
% ============================================================================
\section{Lösungen}

\subsection{AB-06b Lösungen}

\textbf{Merkkasten ir:}
\begin{itemize}
    \item yo $\rightarrow$ voy
    \item tú $\rightarrow$ vas
    \item él/ella $\rightarrow$ va
    \item nosotros $\rightarrow$ vamos
    \item vosotros $\rightarrow$ vais
    \item ellos $\rightarrow$ van
\end{itemize}

\textbf{Merkkasten estar:}
\begin{itemize}
    \item yo $\rightarrow$ estoy
    \item tú $\rightarrow$ estás
    \item él/ella $\rightarrow$ está
    \item nosotros $\rightarrow$ estamos
    \item vosotros $\rightarrow$ estáis
    \item ellos $\rightarrow$ están
\end{itemize}

\textbf{Aufgabe 1: Konjugation ir (4P)}
\begin{enumerate}[label=\alph*)]
    \item Yo \textbf{voy} al cine.
    \item Tú \textbf{vas} a la plaza.
    \item María \textbf{va} al parque.
    \item Nosotros \textbf{vamos} a la biblioteca.
\end{enumerate}

\textbf{Aufgabe 2: Kontraktion al (4P)}
\begin{enumerate}[label=\alph*)]
    \item Voy \textbf{al} supermercado. (a + el)
    \item Vas \textbf{a la} plaza. (kein al)
    \item Va \textbf{al} parque. (a + el)
    \item Vamos \textbf{a la} biblioteca. (kein al)
\end{enumerate}

\textbf{Aufgabe 3: estar + Ort (6P)}
\begin{enumerate}[label=\alph*)]
    \item El cine \textbf{está} en la plaza.
    \item La biblioteca \textbf{está} cerca del parque.
    \item ¿Dónde \textbf{estás}? -- \textbf{Estoy} en casa.
\end{enumerate}

\textbf{Aufgabe 4: Dialog (6P)}
\begin{itemize}
    \item A: ¿Adónde vas?
    \item B: Voy al cine. ¿Y tú?
    \item A: Voy a la plaza. ¿Dónde está el cine?
    \item B: El cine está en la calle Mayor.
    \item A: ¡Gracias! ¡Hasta luego!
    \item B: ¡Adiós!
\end{itemize}

% ============================================================================
% 10. ANLAGEN
% ============================================================================
\section{Anlagen}

\begin{enumerate}
    \item Arbeitsblatt (AB-06b.pdf)
    \item Musterlösung (ML-06b.pdf)
    \item Lehrerhinweise (LH-06b.pdf)
    \item Lernpfad (LP-06b.html)
    \item Stadtplan (Kopiervorlage)
\end{enumerate}

% ============================================================================
% 11. DIDAKTIK-SCORE
% ============================================================================
\newpage
\section{Didaktik-Score}

\begin{scorebox}
\textbf{Bewertung der didaktischen Qualität nach 10 Dimensionen}\\
\textbf{Ziel-Score: $\geq$ 9.0 (Premium-Qualität)}

\vspace{1em}
\begin{tabular}{|c|l|c|c|p{4cm}|}
\hline
\textbf{\#} & \textbf{Dimension} & \textbf{Gew.} & \textbf{Score} & \textbf{Notiz} \\
\hline
1 & Differenzierung & 15\% & 9/10 & [B]/[S]/[E] qualitativ \\
\hline
2 & Taxonomie (AFB) & 10\% & 9/10 & 25/50/25 erfüllt \\
\hline
3 & Lernziel-Alignment & 10\% & 10/10 & RP MV explizit \\
\hline
4 & Aktivierung & 10\% & 9/10 & Entdeckend, Stadtplan \\
\hline
5 & Scaffolding & 10\% & 9/10 & Farbcodierung, gestuft \\
\hline
6 & Feedback-Qualität & 10\% & 9/10 & LP: elaboriert \\
\hline
7 & Sprachliche Klarheit & 10\% & 10/10 & Kl. 7 angemessen \\
\hline
8 & Motivation & 10\% & 9/10 & Alltagsbezug Stadt \\
\hline
9 & Inklusion (UDL) & 10\% & 9/10 & Mehrere Zugänge \\
\hline
10 & Praktikabilität & 5\% & 9/10 & 90 Min. realistisch \\
\hline
\multicolumn{3}{|r|}{\textbf{GESAMT:}} & \textbf{9.1/10} & \\
\hline
\end{tabular}

\vspace{1em}
$\checkmark$ \textcolor{successgreen}{\textbf{FREIGABE} (Score $\geq$ 9.0)}
\end{scorebox}

\end{document}
